% ============================================================
%  CHAPTER 5 — Conclusions
%  Typical length: 4–8 pages
%  Cover: summary of findings, contributions, limitations,
%         and directions for future work.
% ============================================================

% Section are made as examples to show how to write — you have to make sections/subsections as per your need.


\chapter{Conclusions}
\label{chap:conclusions}

[Section are made as examples — make your own sections/subsections headings as per your need.]

% ---- 5.1 Summary of Findings -------------------------------
\section{Summary of Findings}
\label{sec:summary}

% TODO: Concisely summarise the main results from Chapter 4.
% Map each finding back to the objectives listed in Chapter 1.
% Use the past tense ("This thesis derived...", "It was shown that...").

[Summarise your key findings here, linked to each objective.]

% ---- 5.2 Contributions -------------------------------------
\section{Contributions}
\label{sec:contributions}

% TODO: State what is new about your work.
% What does your thesis add to the existing body of knowledge?

[List the original contributions of your thesis.]

% ---- 5.3 Limitations ---------------------------------------
\section{Limitations}
\label{sec:limitations}

% TODO: Honestly describe the assumptions and limitations
% of your work. This shows scientific maturity.

[Describe the limitations of your study here.]

% ---- 5.4 Future Work ---------------------------------------
\section{Future Work}
\label{sec:future_work}

% TODO: Suggest specific directions that follow naturally
% from the limitations or open questions in your work.

[Outline directions for future research here.]


% Uncomment the below lines (if you have used AI for the contents of this chapter) to include an AI use disclosure at the end of the chapter. You can modify the wording as needed, but be transparent about how you used AI tools in writing this chapter.

% \vspace{1.0cm}
% \noindent\rule{\textwidth}{0.4pt}
% \vspace{0.2cm}

% \noindent\textit{Note on AI assistance:} An AI language model was used in writing this chapter as a collaborative tool for [describe how — e.g., drafting, editing, structuring conclusions]. All conclusions drawn are the author's own.
